% 第五章 安全性分析
\section{安全性分析}

\subsection{安全模型}

本系统的安全模型基于以下假设：

\begin{enumerate}[label=H\arabic*:]
    \item \textbf{离散对数困难性}：在 secp256k1 椭圆曲线上，计算离散对数 $\log_G H$ 是困难的。
    \item \textbf{诚实监护人}：监护人不会主动泄露自己的份额 $v_i$ 和盲因子 $r_i$。
    \item \textbf{诚实多数}：在 $t$-of-$n$ 配置中，至少 $t$ 个监护人中存在至少 $t$ 个诚实参与者。
\end{enumerate}

\subsection{安全性声明}

\textbf{声明 1（主密钥机密性）}：在假设 H1 和 H2 下，server 无法获取主密钥 $s$ 的任何信息。

\textbf{证明 sketch}：主密钥 $s$ 在创建后被立即销毁。Server 仅持有加密资产 $E_{\text{AES}}(\text{assets}, s)$（无密钥无法解密）和公开承诺 $C_i = r_iG + v_iH$。要从 $C_i$ 中恢复 $v_i$ 需要计算 $\log_G H$（H1）。同时，$s$ 无法从 $\{v_i\}_{i=1}^{t-1}$ 的信息论安全恢复（Shamir 安全性）。$\square$

\textbf{声明 2（承诺绑定性与刷新正确性）}：在假设 H1 下，监护人无法提交与原有承诺不同的值；刷新后的份额仍对应同一个主密钥。

\textbf{证明 sketch}：绑定等价于 ECDLP 困难性。刷新的正确性由零和多项式的性质保证：$Z_v(0) = Z_r(0) = 0$ 确保主密钥不变；$\Sigma(\lambda_j \cdot C'_j) = C_{\text{master}}$ 由同态加法保证。$\square$

\textbf{声明 3（胁迫码不可区分性）}：攻击者无法判断监护人的两份凭证中哪一份是正常份额、哪一份是胁迫份额。

\textbf{证明 sketch}：正常份额 $(v, r)$ 对应承诺 $C = rG + vH$，胁迫份额 $(v', r')$ 对应承诺 $C' = r'G + v'H$。两者使用相同的 Pedersen 承诺结构，承诺点均随机分布于椭圆曲线上。给定 $C$ 和 $C'$，攻击者无法区分哪一个是正常承诺（参考 Pedersen 承诺的完美隐藏属性）。$\square$

\subsection{威胁模型分析}

\begin{table}[H]
    \centering
    \caption{威胁模型与防御措施}
    \label{tab:threats}
    \begin{tabularx}{\textwidth}{lXc}
        \toprule
        \textbf{攻击类型} & \textbf{描述} & \textbf{防御措施} \\
        \midrule
        Server 入侵 & 攻击者获取 server 全部数据 & AES 加密 + 无密钥存储 + 承诺仅含 EC 点 \\
        \midrule
        份额窃取 & 攻击者窃取一份或多份份额邮件 & 信息论安全：$< t$ 份无法恢复 \\
        \midrule
        胁迫攻击 & 攻击者强迫监护人提交份额 & 胁迫码触发警报并终止继承 \\
        \midrule
        中介攻击 & 攻击者截获份额传输 & ECIES 非对称加密传输 \\
        \midrule
        刷新攻击 & 攻击者伪造刷新增量 & 零和多项式验证 + 承诺一致性检查 \\
        \midrule
        重放攻击 & 攻击者重复使用旧份额 & 每次刷新后旧承诺失效，Pedersen 验证 \\
        \bottomrule
    \end{tabularx}
\end{table}

\subsection{与现有方案的对比}

\begin{table}[H]
    \centering
    \caption{与现有数字资产继承方案的对比}
    \label{tab:comparison}
    \begin{tabularx}{\textwidth}{lXXXX}
        \toprule
        \textbf{维度} & \textbf{本方案} & \textbf{Ledger Recover} & \textbf{公证遗嘱} & \textbf{Safe\{Wallet\}} \\
        \midrule
        信任模型 & 数学（密码学） & 中心化服务 + KYC & 公证机构 & 智能合约 \\
        中心化程度 & 去中心化 & 中心化 & 中心化 & 链上多签 \\
        隐私保护 & 全流程加密 & KYC 可见 & 完全暴露 & 链上公开 \\
        资产范围 & 加密货币+云+文件 & 仅 Ledger 设备 & 通用 & 仅 EVM \\
        抗胁迫 & 胁迫码+警报 & 无 & 无 & 无 \\
        份额刷新 & 零信任同态刷新 & 不可刷新 & — & 不可刷新 \\
        成本 & 无订阅费 & \$9.99/月 & 数千-数万元 & Gas 费 \\
        开源可审计 & 全部开源 & 闭源 & — & 部分开源 \\
        \bottomrule
    \end{tabularx}
\end{table}

本方案在信任模型（数学去信任 vs 中心化信任）、隐私保护（全流程加密 vs 暴露）、抗胁迫能力和资产覆盖范围方面均优于现有方案。
